% Build in this directory: pdflatex -interaction=nonstopmode -halt-on-error class10_handout.tex
\documentclass[a4paper,10pt,pdflatex,ja=standard,jafont=ipaex-type1,
  magstyle=real,scale=1]{bxjsarticle}
\usepackage{lmodern}
\usepackage{amsmath,amssymb,mathtools}
\usepackage{geometry}
\geometry{reset,left=22mm,right=22mm,top=20mm,bottom=17mm,
  headheight=13pt,headsep=6mm,footskip=10mm}
\usepackage{xcolor}
\usepackage{enumitem}
\usepackage{array,booktabs,tabularx}
\usepackage{fancyhdr}

\definecolor{ink}{HTML}{193A5C}
\definecolor{accent}{HTML}{237D82}
\definecolor{soft}{HTML}{EAF3F3}
\definecolor{muted}{HTML}{52616B}
\newcommand{\coursename}{解析学 II}
\newcommand{\handoutnumber}{10}
\pagestyle{fancy}
\fancyhf{}
\fancyhead[L]{\small\color{ink}\coursename{}・授業要約と演習}
\fancyhead[R]{\small\color{ink}第\handoutnumber 回}
\fancyfoot[C]{\small\color{ink}\thepage/2}
\renewcommand{\headrulewidth}{0.3pt}
\setlength{\parindent}{0pt}
\setlength{\parskip}{2pt}
\setlength{\abovedisplayskip}{5pt}
\setlength{\belowdisplayskip}{5pt}
\setlength{\abovedisplayshortskip}{3pt}
\setlength{\belowdisplayshortskip}{4pt}
\setlist[itemize]{leftmargin=1.8em,itemsep=1pt,topsep=2pt,parsep=0pt}

\newcommand{\handouttitle}[1]{%
  \begin{center}
    {\Large\color{ink}#1}\par\vspace{3mm}
    {\color{accent}\rule{0.88\linewidth}{0.7pt}}
  \end{center}\vspace{-2mm}}
\newcommand{\handout}[2]{%
  \renewcommand{\handoutnumber}{#1}%
  \handouttitle{第#1回\quad #2}}
\newcommand{\topic}[1]{\par\vspace{5pt}%
  {\large\sffamily\mdseries\color{accent}#1}\par\vspace{2pt}}
\newcommand{\keybox}[1]{\par\vspace{3pt}\begingroup
  \setlength{\fboxsep}{4pt}%
  \colorbox{soft}{\parbox{\dimexpr\linewidth-2\fboxsep\relax}{\small #1}}%
  \endgroup\par\vspace{2pt}}
\newenvironment{problems}
  {\begin{enumerate}[leftmargin=2em,itemsep=3pt,topsep=3pt,parsep=0pt,
    font=\color{accent}]}
  {\end{enumerate}}


\newcommand{\examplelabel}[1]{{\sffamily\mdseries\color{ink}#1}}
\newcommand{\answerroom}{\par\vspace{16mm}}

\begin{document}
\fontsize{10pt}{13.5pt}\selectfont
\setlength{\abovedisplayskip}{4pt}
\setlength{\belowdisplayskip}{4pt}
\handout{10}{スカラー場とベクトル場}
\topic{場と微分演算子}
開集合 $U\subset\mathbb R^3$ 上の関数 $f:U\to\mathbb R$ をスカラー場、
$\vec F=(P,Q,R):U\to\mathbb R^3$ をベクトル場という。
以下、成分は $C^1$ 級とする。温度・密度はスカラー場、速度・力はベクトル場の例である。

\topic{勾配：スカラー場からベクトル場へ}
$$
 \nabla f=(f_x,f_y,f_z),\qquad
 D_{\vec u}f=\nabla f\cdot\vec u\quad(|\vec u|=1).
$$
方向微分の最大値は $|\nabla f|$。$\nabla f\ne0$ なら、最大になる方向は
$\vec u=\nabla f/|\nabla f|$ である。等位面 $f=c$ の正則点では $\nabla f$ は法線になる。
曲線 $\vec r(t)$ に沿う変化率は連鎖律により
$$
 \frac d{dt}f(\vec r(t))=\nabla f(\vec r(t))\cdot\vec r'(t).
$$

\topic{発散と回転}
発散はスカラー場、回転はベクトル場を与える。
$$
 \operatorname{div}\vec F=\nabla\cdot\vec F=P_x+Q_y+R_z,
$$
$$
 \operatorname{curl}\vec F=\nabla\times\vec F
 =(R_y-Q_z,\ P_z-R_x,\ Q_x-P_y).
$$
速度場の発散は局所的な湧き出し・吸い込み、回転は局所的な渦の向きと強さを表す。
\par\examplelabel{例：$\vec F=(-y,x,0)$。}
$$
 \operatorname{div}\vec F=0,\qquad
 \operatorname{curl}\vec F=(0,0,2).
$$
発散が0でも、回転が0とは限らない。

\topic{ポテンシャルと恒等式}
$\vec F=\nabla\phi$ と書けるとき、$\phi$ をスカラーポテンシャルという。
$C^2$ 級の $\phi$、$\vec A$ について、混合偏微分の交換から
$$
 \nabla\times(\nabla\phi)=\vec0,\qquad
 \nabla\cdot(\nabla\times\vec A)=0.
$$
$\vec F=\nabla\times\vec A$ と書けるとき、$\vec A$ はベクトルポテンシャルである。
星形の開集合上では、$C^1$ 級の $\vec F$ に対する $\operatorname{curl}\vec F=\vec0$ は
スカラーポテンシャルの存在を導く。領域に穴がある場合は、この条件だけでは十分でない。

\topic{例：ポテンシャルを成分から求める}
$\vec F=(2xy,x^2,0)$ に対して $\phi_x=2xy$ を積分すると
$\phi=x^2y+h(y,z)$。残りの成分から $h_y=h_z=0$ なので
$$
 \phi=x^2y+C,\qquad \nabla\phi=(2xy,x^2,0).
$$
\keybox{勾配・発散・回転の入力と出力を確認する。ポテンシャルの存在を判断するときは、場の微分だけでなく領域の条件も調べる。}

\newpage
\handouttitle{第10回\quad 演習問題}
計算過程を記し、必要な条件・積分領域・向きを明示すること。
\begin{problems}
\item $f=x^2+2y^2+3z^2$ の勾配を求めよ。点 $(1,1,1)$ における最大の方向微分と、その方向の単位ベクトルを求めよ。\answerroom
\item $\vec F=(xy,yz,zx)$ の発散と回転を計算し、それぞれがスカラーかベクトルかを確認せよ。\answerroom
\item $\vec F=(yz,xz,xy)$ の発散・回転を求め、スカラーポテンシャルを一つ求めよ。\answerroom
\item $f=x^2-y^2+z^2$ の $\nabla f=\vec0$ となる点を求めよ。その点が極大・極小か、座標軸上での関数値を比較して判断せよ。\answerroom
\item $C^2$ 級の $f$ に対して $\nabla\times(\nabla f)=\vec0$ を、成分ごとに混合偏微分を用いて示せ。\answerroom
\item $\vec F=(-y,x,0)$ のベクトルポテンシャルを、$\vec A=(0,0,h(x,y))$ の形で一つ求めよ。また、この $\vec F$ がスカラーポテンシャルをもたない理由を述べよ。\answerroom
\end{problems}
\end{document}
